\documentclass[11pt,twoside]{article}
\usepackage{etex}

\raggedbottom

%geometry (sets margin) and other useful packages
\usepackage{geometry}
\geometry{top=1in, left=1in,right=1in,bottom=1in}
 \usepackage{graphicx,booktabs,calc}
 
\usepackage{listings}


% Marginpar width
%Marginpar width
\newcommand{\pts}[1]{\marginpar{ \small\hspace{0pt} \textit{[#1]} } } 
\setlength{\marginparwidth}{.5in}
%\reversemarginpar
%\setlength{\marginparsep}{.02in}

 
%\usepackage{cmbright}lstinputlisting
%\usepackage[T1]{pbsi}


\usepackage{chngcntr,mathtools}
\counterwithin{figure}{section}
\numberwithin{equation}{section}

%\usepackage{listings}

%AMS-TeX packages
\usepackage{amssymb,amsmath,amsthm} 
\usepackage{bm}
\usepackage[mathscr]{eucal}
\usepackage{colortbl}
\usepackage{color}


\usepackage{subfigure,hyperref,enumerate,polynom,polynomial}
\usepackage{multirow,minitoc,fancybox,array,multicol}

\definecolor{slblue}{rgb}{0,.3,.62}
\hypersetup{
    colorlinks,%
    citecolor=blue,%
    filecolor=blue,%
    linkcolor=blue,
    urlcolor=slblue
}

%%%TIKZ
\usepackage{tikz}

\usepackage{pgfplots}
\pgfplotsset{compat=newest}

\usetikzlibrary{arrows,shapes,positioning}
\usetikzlibrary{decorations.markings}
\usetikzlibrary{shadows}
\usetikzlibrary{patterns}
%\usetikzlibrary{circuits.ee.IEC}
\usetikzlibrary{decorations.text}
% For Sagnac Picture
\usetikzlibrary{%
    decorations.pathreplacing,%
    decorations.pathmorphing%
}

\tikzstyle arrowstyle=[black,scale=2]
\tikzstyle directed=[postaction={decorate,decoration={markings,
    mark=at position .65 with {\arrow[arrowstyle]{stealth}}}}]
\tikzstyle reverse directed=[postaction={decorate,decoration={markings,
    mark=at position .65 with {\arrowreversed[arrowstyle]{stealth};}}}]
\tikzstyle dir=[postaction={decorate,decoration={markings,
    mark=at position .98 with {\arrow[arrowstyle]{latex}}}}]
\tikzstyle rev dir=[postaction={decorate,decoration={markings,
    mark=at position .98 with {\arrowreversed[arrowstyle]{latex};}}}]

\usepackage{ctable}

%
%Redefining sections as problems
%
\makeatletter
\newenvironment{exercise}{\@startsection 
	{section}
	{1}
	{-.2em}
	{-3.5ex plus -1ex minus -.2ex}
    	{1.3ex plus .2ex}
    	{\pagebreak[3]%forces pagebreak when space is small; use \eject for better results
	\large\bf\noindent{Exercise 1.\hspace{-1.5ex} }
	}
	}
	%{\vspace{1ex}\begin{center} \rule{0.3\linewidth}{.3pt}\end{center}}
	%\begin{center}\large\bf \ldots\ldots\ldots\end{center}}
\makeatother

%
%Fancy-header package to modify header/page numbering 
%
\usepackage{fancyhdr}
\pagestyle{fancy}
%\addtolength{\headwidth}{\marginparsep} %these change header-rule width
%\addtolength{\headwidth}{\marginparwidth}
%\fancyheadoffset{30pt}
%\fancyfootoffset{30pt}
\fancyhead[LO,RE]{\small Oke}
\fancyhead[RO,LE]{\small Page \thepage} 
\fancyfoot[RO,LE]{\small PS 3} 
\fancyfoot[LO,RE]{\small \scshape CEE 260/MIE 273} 
\cfoot{} 
\renewcommand{\headrulewidth}{0.1pt} 
\renewcommand{\footrulewidth}{0.1pt}
%\setlength\voffset{-0.25in}
%\setlength\textheight{648pt}


\usepackage{paralist}

\newcommand{\osn}{\oldstylenums}
\newcommand{\lt}{\left}
\newcommand{\rt}{\right}
\newcommand{\pt}{\phantom}
\newcommand{\tf}{\therefore}
\newcommand{\?}{\stackrel{?}{=}}
\newcommand{\fr}{\frac}
\newcommand{\dfr}{\dfrac}
\newcommand{\ul}{\underline}
\newcommand{\tn}{\tabularnewline}
\newcommand{\nl}{\newline}
\newcommand\relph[1]{\mathrel{\phantom{#1}}}
\newcommand{\cm}{\checkmark}
\newcommand{\ol}{\overline}
\newcommand{\rd}{\color{red}}
\newcommand{\bl}{\color{blue}}
\newcommand{\pl}{\color{purple}}
\newcommand{\og}{\color{orange!90!black}}
\newcommand{\gr}{\color{green!40!black}}
\newcommand{\nin}{\noindent}
\newcommand{\la}{\lambda}
\renewcommand{\th}{\theta}
\newcommand*\circled[1]{\tikz[baseline=(char.base)]{
            \node[shape=circle,draw,thick,inner sep=1pt] (char) {\small #1};}}

\newcommand{\bc}{\begin{compactenum}[\quad--]}
\newcommand{\ec}{\end{compactenum}}

\newcommand{\n}{\\[2mm]}
%% GREEK LETTERS
\newcommand{\al}{\alpha}
\newcommand{\gam}{\gamma}
\newcommand{\eps}{\epsilon}
\newcommand{\sig}{\sigma}

\newcommand{\p}{\partial}
\newcommand{\pd}[2]{\frac{\partial{#1}}{\partial{#2}}}
\newcommand{\dpd}[2]{\dfrac{\partial{#1}}{\partial{#2}}}
\newcommand{\pdd}[2]{\frac{\partial^2{#1}}{\partial{#2}^2}}
\newcommand{\mr}{\mathbb{R}}
\newcommand{\xs}{x^{*}}
\newenvironment{solution}
{\medskip\par\quad\quad\begin{minipage}[c]{.8\textwidth}\gr}{\medskip\end{minipage}}

 
%%%%%%%%%%%%%%%%%%%%%%%%%%%%%%%%%%%%%%%%%%%%%%%%%%%
%%%%%%%%%%%%%%%%%%%%%%%%%%%%%%%%%%%%%%%%%%%%%%%%%%%

\begin{document}

\lstset{language=C++,
                basicstyle=\tiny\ttfamily,
                keywordstyle=\color{blue}\ttfamily,
                stringstyle=\color{red}\ttfamily,
                commentstyle=\color{gray}\ttfamily,
                morecomment=[l][\color{gray}]{\#}
}


\thispagestyle{empty}


\nin{\LARGE Problem Set 3 {\gr }}\hfill{\bf Prof. Oke}

\medskip\hrule\medskip

\nin {\small CEE 260/MIE 273: Probability \& Statistics in Civil Engineering
\hfill\textit{ 09.24.2026}}

\bigskip

\nin{\it Due Wednesday, September 30, 2026 at 11:59 PM as PDF uploaded via Gradescope.
  If it helps and if possible, you can write your responses directly on this document and upload it instead.
    \textbf{Show as much work as possible in order to get FULL credit.}
  There are 4 problems with a total of 32 points available.
}\\


\section*{Problem 1 \textit{(8 points)}}

Respond ``T'' ({\it True})  or  ``F'' (\textit{False}) to the following statements. Use the boxes provided. Each response is worth 1 point.

\begin{enumerate}[\bf (i)]
\item \hfill
  \begin{minipage}{.1\linewidth}
    \framebox(40,40){ \gr  }
  \end{minipage}\quad
  \begin{minipage}{.85\linewidth}
    Given three events $D$, $A$ and $G$. If $P(\ol{D}|AG) = 1$, then $P(D)$ is an impossible event.
 
  \end{minipage}
%  \underline{\hspace{10ex}} %F 

    \smallskip
  
\item \hfill
  \begin{minipage}{.1\linewidth}
    \framebox(40,40){\gr  }
  \end{minipage}\quad
  \begin{minipage}{.85\linewidth}
   Two events $A$ and $B$ can both be mutually exclusive and yet collectively exhaustive.
  \end{minipage}

  \smallskip
  
\item \hfill
  \begin{minipage}{.1\linewidth}
    \framebox(40,40){\gr  }
  \end{minipage}\quad
  \begin{minipage}{.85\linewidth}
    If $P(A|B) = P(A)$ for a set of events $A$ and $B$, then both events are  dependent.
  \end{minipage}

  \smallskip
  
\item \hfill
  \begin{minipage}{.1\linewidth}
    \framebox(40,40){\gr  }
  \end{minipage}\quad
  \begin{minipage}{.85\linewidth}
    Events $E$ and $F$ are mutually exclusive and collectively exhaustive. If $P(E) = 0.3$, then $P(F) = 0.7$.
   \end{minipage}

  \smallskip

  \item \hfill
  \begin{minipage}{.1\linewidth}
    \framebox(40,40){\gr  }
  \end{minipage}\quad
  \begin{minipage}{.85\linewidth}
    Events $E$ and $F$ are  collectively exhaustive but \textit{not} mutually exclusive.
    For a given event $A$, its probability can be obtained by $P(A) = P(AE) + P(AF)$.
    % {\gr This only holds if $E$ and $F$ are mutually exclusive. But since they are not, then $P(A) = P(AE\cup AF) =
    %   P(AE) + P(AF) - P(AE\cap AF)$ (addition rule).}
  \end{minipage}
  
  \smallskip
  \item \hfill
  \begin{minipage}{.1\linewidth}
    \framebox(40,40){\gr  }
  \end{minipage}\quad
  \begin{minipage}{.85\linewidth}
    The number of ways 6 books can be arranged on a bookshelf is 720. If two of the books are identical, then the total number of distinct arrangements is 360. %{\gr The permutations are $\fr{n!}{n_{1}!}$, where $n_{1} = 2$. Thus $\fr{6!}{2} = 360$.}
  \end{minipage}
  
  \smallskip
  
\item \hfill
  \begin{minipage}{.1\linewidth}
    \framebox(40,40){\gr  }
  \end{minipage}\quad
  \begin{minipage}{.85\linewidth}
    The number of distinct subgroups of size $n$ that can be formed from a larger group of $m$ objects is given by $\fr{m!}{n!(n-m)!}$.
    %{\gr (The correct answer is $\fr{m!}{n!(m-n)!}$)}
  \end{minipage}

  \smallskip

  \item \hfill
  \begin{minipage}{.1\linewidth}
    \framebox(40,40){\gr  }
  \end{minipage}\quad
  \begin{minipage}{.85\linewidth}
    The events $E_{1}$ and $E_{2}$ are independent. If $P(E_{1}) = 0.4$ and $P(E_{1}E_{2}) = 0.04$, then $P(E_{2}) = 0.1$.
    %{\gr ($P(E_{2}) = P(E_{1}E_{2})/P(E_{1}) = 0.04/0.4 = 0.1$)}
  \end{minipage}
  \end{enumerate}
% \item \hfill
%   \begin{minipage}{.1\linewidth}
%     \framebox(40,40){\gr T}   The airline industry in a certain country is subject to labor strikes by the

%   \end{minipage}\quad
%   \begin{minipage}{.85\linewidth}
%     The number of ways 6 books can be arranged on a bookshelf is 720. If two of the books are identical, then the total number of distinct arrangements is 360. {\gr The permutations are $\fr{n!}{n_{1}!}$, where $n_{1} = 2$. Thus $\fr{6!}{2} = 360$.}
%   \end{minipage}
  
  %\smallskip
  
% \item \hfill
%   \begin{minipage}{.1\linewidth}
%     \framebox(40,40){\gr  }
%   \end{minipage}\quad
%   \begin{minipage}{.85\linewidth}
%     The number of tails $X$ in 50 tosses of a coin is a continuous random variable.
%   \end{minipage}

%   \smallskip

%   \item \hfill
%   \begin{minipage}{.1\linewidth}
%     \framebox(40,40){\gr  }
%   \end{minipage}\quad
%   \begin{minipage}{.85\linewidth}
%   The length of time $T$ between car accidents at a busy intersection is a continuous random variable.
%   \end{minipage}
%   \end{enumerate}

  % \smallskip
  
% \item \hfill
%   \begin{minipage}{.1\linewidth}
%     \framebox(40,40){}
%   \end{minipage}\quad
%   \begin{minipage}{.85\linewidth}
%     A good clustering solution should maximize the total within-cluster variation. %False
%   \end{minipage}

  \eject

   \section*{Problem 2 \textit{(9 points)}}
   The airline industry in a certain country is subject to labor strikes by the
    pilots (event $A$), mechanics (event $B$) or flight attendants (event $C$).
    The probability of strikes by each of the individual groups in the next 3 years is given by:
    \[P(A) = 0.03\quad P(B) = 0.05\quad P(C) = 0.06\]
    Assuming that events $A$, $B$ and $C$ are statistically independent, find the probability that:

    \begin{enumerate}[\bf (a)]
    \item only pilots will strike in the next 3 years \pts{1} \vspace{25ex}
    \item both mechanics and flight attendants will strike in the next 3 years \pts{3} \vspace{25ex}
    \item any 2 of the 3 groups will strike in the next 3 years \pts{5} \vspace{25ex}
    \end{enumerate}
 


% \section*{Problem 2 \textit{(17 points)}}
% For a construction project, you are given the following conditional
% probabilities: $P(E|G)$, $P(E|N)$, $P(E|BC)$ and $P(E|B\ol{C})$, where $E$
% denotes project completion, $G$ good weather, $N$ normal weather $B$ bad
% weather, and $C$ the launching of a crash program to speed up completion.
% \begin{align*}
%   P(E|G) &= 1 \\
%   P(E|N) &= 0.9 \\
%   P(E|BC) &=  0.8 \\
%   P(E|B\ol{C}) &= 0.2 \\
%   P(G) & = 0.2 \\
%   P(N) &= 0.4 \\
%   P(B) &= 0.4 \\
%   P(C) &= 0.5 
% \end{align*}

% \begin{enumerate}[(a)]
% \item How would you chracterize the event $(E|G)$? \pts{1}
%   \vspace{5ex}

% \item How would characterize the events $E|BC$ and $E|B\ol{C}$? \pts{1}
%   \vspace{5ex}

% \item $P(G) + P(N) + P(B) = 1$. How would you characterize the 3 events $G$, $N$ and $B$? \pts{1}
%   \vspace{5ex}

% \item What is the probability of the complement of $C$, i.e.\ $\ol{C}$? \pts{1}
%   \vspace{5ex}

% \item How would you characterize $C$ and $\ol{C}$? \pts{1}
%   \vspace{5ex}

% \item The Venn diagram below represents all the events in this problem. \pts{5} Match the description of each region to the appropriate event.

%        \begin{tikzpicture}[scale=.8]
%     \draw[thick] (0,0) rectangle (10,4) node[above left] {$\bm S$};
% %    \draw[thick] (1.5,0) ellipse  (4 cm and 2 cm) node[] {$\bm{D}$};
%     \filldraw[thick, fill=orange!30!white] (0,0) -- (2,0) -- (2,4) -- (0,4) -- cycle;
%     \filldraw[thick, fill=purple!20!white] (2,0) -- (6,0) -- (6,4) -- (2,4) -- cycle;
%     \filldraw[thick, fill=green!20!white] (6,0) -- (10,0) -- (10,4) -- (6,4) -- cycle;
%     \filldraw[red, thick,pattern = horizontal lines, pattern color=red] (0,0) -- (10,0) -- (10,2) -- (0,2) -- cycle;    
%     \filldraw[blue, thick,pattern=north west lines,pattern color=blue!50!white] (0,0) -- (2,0) -- (9,2) -- (2,4) -- (0,4) -- cycle;
%     % \node[blue] at (3, 2.5) {$\bm E$};
%     % \node[orange] at (1,3.75) {$\bm{G}$};
%     % \node[purple] at (4,3.75) {$\bm{N}$};
%     % \node at (8,3.75) {$\bm{B}$};
%     % \node[red] at (7,1) {$\bm{C}$};
    
%     %\begin{scope}
%     %\clip   (1.5,0) ellipse  (4 cm and 2 cm);
%     %\fill[pattern=north west lines, pattern color=green!50!black] (4,-3) rectangle (7,3);
%     %\end{scope}
% %    \draw[thick] (3,0) circle (2 cm) node[right] {$\bm{E_2}$};
%   \end{tikzpicture}

%   \begin{center}
%     \begin{tabular}{p{5.5in} p{.5in}}
%       \fbox{Orange rectangle to the left (20\% of  sample space)} & \fbox{\bf  B} \\[3mm]
%        \fbox{Purple rectangle in the middle (40\% of sample space)} &  \fbox{\bf C} \\[3mm]
%       \fbox{Green rectangle on the right (40\% of sample space)} & \fbox{\bf  E} \\[3mm]
%        \fbox{Blue pentagon hatched with diagonal lines} &  \fbox{\bf G} \\[3mm]
%        \fbox{Red rectangle hatched with horizontal lines (50\% of sample space)} &  \fbox{\bf N} 
                                                             
%     \end{tabular}
%   \end{center}
%   % \begin{itemize}
%   % \item $G$ is the orange rectangle to the left (20\% of  sample space)
%   % \item $N$ is the purple rectangle in the middle (40\% of sample space)
%   % \item $B$ is the green rectangle on the right (40\% of sample space)
%   % \item $E$ is the blue pentagon hatched with diagonal lines
%   % \item $C$ is the red rectangle hatched with horizontal lines (50\% of sample space)
%   % \end{itemize}

% % \item Draw and label a Venn diagram illustrating the space of events. \pts{3}
% %   \vspace{40ex}
  
% \item Use the theorem of total probability to show that $P(E) = 0.76$.\pts{4}
%   Hints:
%   \begin{itemize}
%   \item $ P(E) = P(E|G)P(G) + P(E|N)P(N) + P(E|BC)P(BC) + P(E|B\ol{C})P(B\ol{C})$
%   \item $B$ and $C$ are independent events.
%   \end{itemize}
%   % \begin{align*}
%   %   P(E) &= P(E|G)P(G) + P(E|N)P(N) + P(E|BC)P(BC) + P(E|B\ol{C})P(B\ol{C}) \\
%   %        &= 1(0.2) + 0.9(0.4) + 0.8(0.4)(0.5) + 0.2(0.4)(0.5) \\
%   %        &= \boxed{0.76}
%   %  \end{align*}
%  \vspace{50ex}
  

% \item Find $P(N|\ol{E})$. \pts{3}
%   % \begin{align*}
%   %   P(N|\ol{E}) &= \fr{P(\ol{E}|N)P(N)}{P(\ol{E})} \\
%   %               &= \fr{(1-0.9)(0.4)}{1-0.76} = \boxed{0.167}
%   % \end{align*}
% \end{enumerate}

\eject

\section*{Problem 3 \textit{Bayes' Theorem (8 points)}}
Given an earthquake of intensity
\begin{equation}
  \label{eq:10}
  X :  \{\text{light } (L), \text{moderate } (M), \text{important } (I)\}
\end{equation}
and a structure that can be in a state
\begin{equation}
  \label{eq:11}
  Y :   \{\text{damaged} (D), \text{undamaged} (\ol{D})\}
\end{equation}
The likelihood of damage given earthquake intensity is given by the following conditional probabilities:
\begin{align*}
  \label{eq:12}
  P(D|L) &= 0.01 \\
  P(D|M) &= 0.10 \\
  P(D|I) &= 0.60 
  \end{align*}
and the prior probability of each intensity is given by
\begin{align*}
    P(L) &= 0.90 \\
    P(M) &= 0.08 \\
    P(I) &= 0.02 
\end{align*}
% An earthquake has just occurred and damaged buildings have been observed, i.e.\
% \begin{equation}
%   \label{eq:14}
%   y = D
% \end{equation}

\begin{enumerate}[\bf (a)]
\item Draw and label a Venn diagram illustrating the given events. \pts{5}
 \vspace{40ex}

\item Find the total probability $P(D)$. \pts{3}
  \vspace{20ex}
\end{enumerate}

\eject
\section*{Problem 4 \textit{Bayes' Theorem (continued; 7 points)}}
Use the quantities provided and the results from Problem to answer the following questions.
\begin{enumerate}[\bf (a)]
\item Use Bayes' Theorem to find the posterior probabilities $P(L|D)$, $P(M|D)$ and $P(I|D)$.\pts{6}
  \vspace{50ex}

\item Show that the probabilities $P(L|D)$, $P(M|D)$ and $P(I|D)$ sum up to 1. \pts{1}
  \vspace{15ex}


\end{enumerate}



% \vfill

    
    
% \end{document}

\end{document}

%%% Local Variables:
%%% mode: latex
%%% TeX-master: t
%%% End:
